\documentclass[11pt]{article}
\title{Nonconsecutive Subset Sums Modulo Three: Exact Encodings and a Candidate Linear Recurrence}
\date{}
\author{SIGMA AI-assisted research synthesis}
\usepackage[T1]{fontenc}
\usepackage[utf8]{inputenc}
\usepackage[margin=1in]{geometry}
\usepackage{amsmath,amssymb,booktabs,longtable}
\usepackage{microtype}
\usepackage{titlesec}
\emergencystretch=2em
\raggedbottom
\linespread{1.07}
\usepackage{xcolor}
\definecolor{blue}{HTML}{24485E}
\titleformat{\section}{\Large\bfseries\sffamily\color{blue}}{\thesection}{0.8em}{}
\titleformat{\subsection}{\large\bfseries\sffamily}{\thesubsection}{0.7em}{}
\setlength{\parskip}{.6em}
\setlength{\parindent}{0pt}
\begin{document}
\maketitle

\begin{abstract}
Let $a_n$ count subsets of \$\textbackslash{}\{1,\textbackslash{}ldots,n\textbackslash{}\}\$ with no consecutive elements and sum divisible by $3$, including the empty subset. Weighted subset polynomials connect this problem to Fibonacci counting and a roots-of-unity filter, while a reversible gap-removal map gives a complementary Gaussian-binomial representation. The evidence supports the exponent correction required by the standard Gaussian-binomial convention and reports the exact finite count $a_{12}=128$. It also rejects the short recurrence $a_{n+3}=a_{n+2}+a_n$ by an explicit counterexample. A more substantial candidate is $a_{n+12}-5a_{n+9}+2a_{n+6}+5a_{n+3}+a_n=0$, supported by bounded exact computations.

The proposed proof through periodic transfer matrices remains incomplete: ordered phase products, boundary cases, and a sufficient initial-value specification still require justification. The requested limit $a_n/F_{n+2}\to 1/3$ likewise remains unresolved because the nontrivial root-of-unity contributions lack a proved asymptotic bound. The planned article therefore presents exact structural reductions and finite results without promoting the recurrence candidate or limiting density to theorems.
\end{abstract}
\section*{Evidence coverage limitation}
The evidence index omitted 996 material saved research field fragment(s) across 6 source(s). These fragments were not reviewed for this paper. Interpretations that could depend on omitted material remain provisional; the complete frozen source snapshot and omission ledger are retained with the generation record.
\section{Introduction}
For each integer $n\geq 0$, count the subsets \$S\textbackslash{}subseteq\textbackslash{}\{1,\textbackslash{}ldots,n\textbackslash{}\}\$ containing no two consecutive integers and satisfying \$\textbackslash{}sum\_\{i\textbackslash{}in S\}i\textbackslash{}equiv 0\textbackslash{}pmod 3\$. The empty subset contributes, so $a_0=1$. The objective is an integer constant-coefficient recurrence with sufficient initial values, an exact determination of $a_{12}$, and a proof that $a_n/F_{n+2}\to 1/3$, where $F_0=0$ and $F_1=1$. [C01]

The central mathematical distinction is between an exact encoding of the counting problem and a proved global statement about its solution. The available evidence contains useful encodings, finite computations, and a falsified early recurrence, but its reviews leave the general recurrence and asymptotic argument incomplete. The article should therefore develop the supported reductions and identify precisely where the proposed proof stops, rather than present a completed solution. [C13]

\section{Evidence and scope}
Use a six-state count indexed by $(r,b)$, where \$r\textbackslash{}in\textbackslash{}\{0,1,2\textbackslash{}\}\$ records the selected-index sum modulo $3$ and \$b\textbackslash{}in\textbackslash{}\{0,1\textbackslash{}\}\$ records whether the preceding integer was selected. Processing the next integer permits omission from either previous state, but inclusion only when $b=0$; inclusion adds that integer to the residue. This supplies an exact finite counting method and a transfer rule whose dependence on the processed integer has period $3$. [C04]

Use weighted subset polynomials and the roots-of-unity filter to separate the total count from the two nontrivial residue modes. The filter averages the evaluations at $1$, $\omega$, and $\omega^2$, with $\omega^3=1$ and \$\textbackslash{}omega\textbackslash{}ne1\$; the evaluation at $1$ is the Fibonacci count. This reduction identifies the components that a recurrence proof and an asymptotic proof must control. [C03]

Use gap removal as a complementary combinatorial representation, retaining the standard partition-normalized Gaussian-binomial convention. The position shift contributes $k(k-1)/2$ to the exponent, while the ordinary subset baseline contributes $k(k+1)/2$, giving the combined exponent $k^2$. This correction must replace the earlier inconsistent exponent explanation. [C06]

Treat the saved finite enumerations, polynomial comparisons, and recurrence residuals as bounded evidence, not as proofs for arbitrary indices. Distinguish a reported exact finite value from a general theorem, and distinguish agreement among encodings from independent verification of the missing matrix and growth arguments. [C13]

\section{Weighted Counting and Residue Extraction}
The weighted subset polynomial gives an exact connection between the adjacency restriction, Fibonacci counting, and extraction of the required residue class. For each integer $n\geq 0$, let $\mathcal A_n$ be the family of subsets of \$\textbackslash{}\{1,\textbackslash{}ldots,n\textbackslash{}\}\$ containing no two consecutive integers, with \$\textbackslash{}mathcal A\_0=\textbackslash{}\{\textbackslash{}varnothing\textbackslash{}\}\$. Write $\sigma(S)=\sum_{i\in S}i$ and define $P_n(q)=\sum_{S\in\mathcal A_n}q^{\sigma(S)}$. The empty subset has sum $0$ and contributes the constant term $1$; in particular, it is included in \$a\_n=\textbackslash{}\#\textbackslash{}\{S\textbackslash{}in\textbackslash{}mathcal A\_n:\textbackslash{}sigma(S)\textbackslash{}equiv 0\textbackslash{}pmod 3\textbackslash{}\}\$.

For $n\geq 2$, partition $\mathcal A_n$ according to whether a subset contains $n$. Subsets omitting $n$ are precisely the members of $\mathcal A_{n-1}$, with unchanged weights. A subset containing $n$ must omit $n-1$, so deleting $n$ leaves a member of $\mathcal A_{n-2}$; conversely, adjoining $n$ to any member of $\mathcal A_{n-2}$ produces an admissible subset. This correspondence increases the sum by $n$, and the two disjoint classes therefore give $P_n(q)=P_{n-1}(q)+q^nP_{n-2}(q)$ for $n\geq 2$.

The initial polynomials must be specified separately: $P_0(q)=1$ and $P_1(q)=1+q$. The first records only the empty subset, while the second records the empty subset and \$\textbackslash{}\{1\textbackslash{}\}\$. Thus the polynomial recurrence requires no negative-index convention and is not being asserted at $n=0$ or $n=1$.

Let $\omega$ be a primitive cube root of unity, so that $\omega^3=1$, \$\textbackslash{}omega\textbackslash{}ne 1\$, and $1+\omega+\omega^2=0$. For every nonnegative integer $m$, the average $(1+\omega^m+\omega^{2m})/3$ equals $1$ when \$m\textbackslash{}equiv 0\textbackslash{}pmod 3\$ and equals $0$ otherwise. Indeed, in the divisible case all three summands are $1$; in either remaining residue class their sum is $1+\omega+\omega^2$.

Applying this identity to the exponent of each subset weight gives the exact filter $a_n=(P_n(1)+P_n(\omega)+P_n(\omega^2))/3$ for every $n\geq 0$. The justification is termwise: a subset contributes one to the average precisely when its sum is divisible by $3$, and otherwise contributes zero. Thus averaging the evaluations selects all exponents in the required residue class, rather than extracting a single coefficient of $P_n(q)$.

At $q=1$, every admissible subset has weight $1$, and the polynomial recurrence becomes $P_n(1)=P_{n-1}(1)+P_{n-2}(1)$, with initial values $P_0(1)=1$ and $P_1(1)=2$. With Fibonacci numbers defined by $F_0=0$, $F_1=1$, and $F_{j+2}=F_{j+1}+F_j$ for $j\geq 0$, these initial values and the recurrence establish $P_n(1)=F_{n+2}$ by induction. Consequently, $a_n=(F_{n+2}+P_n(\omega)+P_n(\omega^2))/3$: the unweighted contribution is determined exactly, and the entire effect of the divisibility condition beyond one third of the total count is isolated in the two nontrivial root-of-unity evaluations.

\section{Gap Removal and the Gaussian-Binomial Weight}
Gap removal gives an exact cardinality-by-cardinality representation of the weighted polynomial $P_n(q)$, separating the adjacency restriction from the normalization of the Gaussian binomial coefficient. The essential point is that two distinct shifts contribute to the weight: removing compulsory gaps changes the subset sum, while the standard Gaussian polynomial removes the minimum sum of an ordinary subset. Keeping these shifts separate explains the exponent $k^2$.

Fix a nonempty admissible subset with ordered elements \$1\textbackslash{}le i\_1\textless{}\textbackslash{}cdots\textless{}i\_k\textbackslash{}le n\$, so that \$i\_\{j+1\}-i\_j\textbackslash{}ge2\$, and define $y_j=i_j-(j-1)$. Then \$y\_\{j+1\}-y\_j=i\_\{j+1\}-i\_j-1\textbackslash{}ge1\$, so the shifted positions are strictly increasing. Their endpoints satisfy \$y\_1=i\_1\textbackslash{}ge1\$ and \$y\_k=i\_k-(k-1)\textbackslash{}le n-k+1\$, placing them in the ordinary subset range \$1\textbackslash{}le y\_1\textless{}\textbackslash{}cdots\textless{}y\_k\textbackslash{}le n-k+1\$.

Conversely, given any ordinary subset \$1\textbackslash{}le y\_1\textless{}\textbackslash{}cdots\textless{}y\_k\textbackslash{}le n-k+1\$, set $i_j=y_j+(j-1)$. Its endpoints satisfy \$i\_1\textbackslash{}ge1\$ and \$i\_k\textbackslash{}le n\$, and its successive gaps satisfy \$i\_\{j+1\}-i\_j=y\_\{j+1\}-y\_j+1\textbackslash{}ge2\$. These formulas undo the forward shift coordinate by coordinate, proving a bijection between admissible subsets of size $k$ and ordinary subsets of size $k$ in \$\textbackslash{}\{1,\textbackslash{}ldots,n-k+1\textbackslash{}\}\$.

The allowable sizes follow from the same spacing condition. Since \$i\_1\textbackslash{}ge1\$ and each subsequent element is at least two larger, \$i\_k\textbackslash{}ge2k-1\$; hence \$2k-1\textbackslash{}le n\$, or \$k\textbackslash{}le\textbackslash{}lceil n/2\textbackslash{}rceil\$. Conversely, whenever \$1\textbackslash{}le k\textbackslash{}le\textbackslash{}lceil n/2\textbackslash{}rceil\$, the subset \$\textbackslash{}\{1,3,\textbackslash{}ldots,2k-1\textbackslash{}\}\$ is admissible, so this upper bound is attained. The empty subset supplies the separate case $k=0$.

To specify the weight convention precisely, use the partition-normalized Gaussian polynomial, characterized for \$0\textbackslash{}le k\textbackslash{}le m\$ by
\$\textbackslash{}displaystyle \{m\textbackslash{}brack k\}\_q=\textbackslash{}sum\_\{1\textbackslash{}le y\_1\textless{}\textbackslash{}cdots\textless{}y\_k\textbackslash{}le m\}q\textasciicircum{}\{\textbackslash{}sum\_\{j=1\}\textasciicircum{}\{k\}y\_j-k(k+1)/2\}.\$
Thus the ordinary subset sum enumerator is \$\textbackslash{}sum\_\{1\textbackslash{}le y\_1\textless{}\textbackslash{}cdots\textless{}y\_k\textbackslash{}le m\}q\textasciicircum{}\{\textbackslash{}sum\_j y\_j\}=q\textasciicircum{}\{k(k+1)/2\}\{m\textbackslash{}brack k\}\_q\$. The baseline is $k(k+1)/2$, the sum of the smallest ordinary subset \$\textbackslash{}\{1,\textbackslash{}ldots,k\textbackslash{}\}\$; it is not $k(k-1)/2$. For $k=0$, both the empty sum and the Gaussian polynomial contribute the constant term $1$.

Under gap removal, the original sum satisfies
$\displaystyle \sum_{j=1}^{k}i_j=\sum_{j=1}^{k}y_j+\sum_{j=1}^{k}(j-1)=\sum_{j=1}^{k}y_j+\frac{k(k-1)}2.$
Consequently, the total weight of admissible subsets of size $k$ is
\$\textbackslash{}displaystyle q\textasciicircum{}\{k(k-1)/2\}\textbackslash{},q\textasciicircum{}\{k(k+1)/2\}\{n-k+1\textbackslash{}brack k\}\_q=q\textasciicircum{}\{k\textasciicircum{}2\}\{n-k+1\textbackslash{}brack k\}\_q.\$
The first factor comes from restoring the compulsory gaps, whereas the second restores the ordinary subset baseline removed by the Gaussian normalization.

Summing over the possible sizes therefore proves the polynomial identity
\$\textbackslash{}displaystyle P\_n(q)=\textbackslash{}sum\_\{k=0\}\textasciicircum{}\{\textbackslash{}lceil n/2\textbackslash{}rceil\}q\textasciicircum{}\{k\textasciicircum{}2\}\{n-k+1\textbackslash{}brack k\}\_q\textbackslash{}qquad(n\textbackslash{}ge0).\$
This is an identity before any specialization of $q$: the reversible map accounts for every admissible subset and preserves its weight with the stated correction. In particular, $k^2$ is the minimum sum $1+3+\cdots+(2k-1)$ for an admissible subset of size $k$. The Gaussian-binomial representation thus supplies a second exact encoding for subsequent residue evaluations, with its exponent fixed by the combinatorics rather than by a normalization-dependent guess.

\section{Exact Finite Counts and a Rejected Recurrence}
The six-state model separates the adjacency restriction from the sum residue, providing an exact counting mechanism against which proposed recurrences can be tested. Its two coordinates retain different information: the residue determines whether a completed subset contributes to $a_n$, whereas the selection bit determines whether the next index may be included. The reported finite counts obtained from this framework support an exact finite answer and decisively reject one short recurrence, without establishing a recurrence for the entire sequence.

Let $D_i(r,b)$ count admissible subsets of \$\textbackslash{}\{1,\textbackslash{}ldots,i\textbackslash{}\}\$ whose sum has residue \$r\textbackslash{}in\textbackslash{}\{0,1,2\textbackslash{}\}\$ modulo $3$, with \$b\textbackslash{}in\textbackslash{}\{0,1\textbackslash{}\}\$ indicating whether $i$ is selected. Before any index is processed, use the initialization $D_0(0,0)=1$, with all other states zero; the bit at this boundary is a convention that permits the first index to be selected. This initialization counts the empty subset exactly once.

When processing index $i\geq1$, omission leaves the residue unchanged and sets the new bit to zero, regardless of the previous bit. Inclusion is permitted only from a state with previous bit zero, and it increases the residue by $i$ modulo $3$. Consequently, with residue arguments interpreted modulo $3$, the transitions are $D_i(r,0)=D_{i-1}(r,0)+D_{i-1}(r,1)$ and $D_i(r,1)=D_{i-1}(r-i,0)$. These cases are disjoint because the resulting subsets differ in whether they contain $i$; they are exhaustive because removing a selected $i$ leaves an admissible subset that omits $i-1$. Thus the state transitions preserve the intended counts, and the desired total is $a_i=D_i(0,0)+D_i(0,1)$.

The reported direct enumeration through $n=8$ gives $(a_0,a_1,a_2,a_3,a_4,a_5,a_6,a_7,a_8)=(1,1,1,2,3,5,8,12,19)$. In that enumeration, a bit mask specifies the selected one-based indices, masks containing adjacent selected positions are rejected, and the remaining masks contribute precisely when their index sum is divisible by $3$. The empty mask is included. The later finite dynamic-programming calculation reports the exact result $a_{12}=128$, also reported from the polynomial representation; these finite computational reports do not by themselves constitute an independently established identity for an infinite sequence.

The early values suggested the tentative rule $a_{n+3}=a_{n+2}+a_n$. It agrees with the first two instances, since $a_3=a_2+a_0=2$ and $a_4=a_3+a_1=3$, but the next instance fails: at $n=2$, the reported counts give $a_5=5$, whereas $a_4+a_2=3+1=4$. This single counterexample refutes the proposed rule as a recurrence valid for all $n\geq0$; agreement at the preceding indices cannot repair that failure.

The negative result changes the role of the finite counts: they are tests that can eliminate a proposed law, not a substitute for deriving one. The six-state transitions retain the structural information that the rejected scalar rule lacked, and their dependence on the processed index is periodic modulo $3$. A replacement recurrence must therefore be justified through a general identity for those transitions, rather than inferred from the initial numerical pattern.

\section{Periodic Transfer and the Order-Twelve Candidate}
Periodic transfer provides a specific algebraic explanation for the candidate recurrence $a_{n+12}-5a_{n+9}+2a_{n+6}+5a_{n+3}+a_n=0$, proposed for every $n\geq 0$. Its largest index displacement is twelve, so this is an order-twelve recurrence in the original index, although every displacement is a multiple of three. The distinction matters: the proposal advances each residue-class subsequence separately, rather than relating adjacent terms of the full sequence. No minimal-order assertion accompanies it.

The proposed mechanism begins with the decomposition $3a_n=P_n(1)+P_n(\omega)+P_n(\omega^2)$, where $\omega$ is a primitive cube root of unity. Here the decomposition is used not to repeat residue extraction, but to identify components that might admit different annihilating relations under three-step advancement. The weighted recurrence $P_n(q)=P_{n-1}(q)+q^nP_{n-2}(q)$, valid for $n\geq 2$, has coefficients periodic with period three at the nontrivial cube roots. This periodicity motivates grouping the transfer operations into blocks of three.

The proposed block polynomials are $x^2-4x-1$ for the unweighted component $P_n(1)=F_{n+2}$ and $x^2-x-1$ for each nontrivial evaluation. These are ingredients of the proof candidate, not established identities for the ordered block matrices. In particular, the asserted trace $1$ and determinant $-1$ of the nontrivial three-step blocks still require symbolic verification for each starting phase. If those assertions hold, the resulting quadratic matrix identity would supply the claimed relation for the corresponding scalar evaluations.

To make the intended combination precise, let $S$ denote three-step advancement, so that $(Su)_n=u_{n+3}$, and let $I$ be the identity operator. The required component relations would be \$(S\textasciicircum{}2-4S-I)P\_\{\textbackslash{}bullet\}(1)=0\$ and \$(S\textasciicircum{}2-S-I)P\_\{\textbackslash{}bullet\}(\textbackslash{}omega\textasciicircum{}j)=0\$ for \$j\textbackslash{}in\textbackslash{}\{1,2\textbackslash{}\}\$. Their algebraic combination is explicit: $(x^2-4x-1)(x^2-x-1)=x^4-5x^3+2x^2+5x+1$. Because polynomials in the same shift operator commute, their product would annihilate every component and hence their average. Substituting $S$ into this product gives exactly the proposed order-twelve recurrence.

The missing bridge is therefore a phase-sensitive transfer identity, not the polynomial multiplication. The transfer matrices do not commute in general, so periodic coefficients alone do not justify replacing successive transfers by an unspecified common block. The ordered product must be stated for each starting residue class and shown to satisfy the appropriate quadratic relation. The initial boundary also needs separate treatment: an argument using a state at $n=0$ must be justified from the initial weighted polynomials, rather than by extending the weighted recurrence below its stated domain.

The recorded exact-integer dynamic program supplies bounded corroboration of the resulting scalar relation. Counts were computed through $n=60$, and the residual $a_{n+12}-5a_{n+9}+2a_{n+6}+5a_{n+3}+a_n$ vanished at every tested index $0\leq n\leq48$. This establishes consistency throughout that range, including the first proposed recurrence instance, but it does not establish the ordered matrix identity or its validity at arbitrary indices.

The initial-data requirement can nevertheless be specified precisely. For each \$r\textbackslash{}in\textbackslash{}\{0,1,2\textbackslash{}\}\$, write $b_m^{(r)}=a_{3m+r}$; the candidate then becomes $b_{m+4}^{(r)}-5b_{m+3}^{(r)}+2b_{m+2}^{(r)}+5b_{m+1}^{(r)}+b_m^{(r)}=0$. Four starting values for each subsequence, equivalently the twelve values from $a_0$ through $a_{11}$, would determine the sequence if the recurrence were established on the proposed domain. The recorded computation gives $(a_0,\ldots,a_{11})=(1,1,1,2,3,5,8,12,19,31,49,79)$; their availability does not replace the outstanding algebraic boundary justification.

Thus the evidence narrows the recurrence problem to a concrete implication: verified quadratic relations for the phase-ordered blocks would yield the displayed integer-coefficient annihilator by a fully specified polynomial calculation. The central advance is an explicit coefficient mechanism and a sufficient initial-data specification, rather than an inference from numerical agreement to an infinite recurrence.

\section{The Remaining Growth Estimate}
The density question reduces to a growth comparison between the nontrivial root-of-unity evaluations and the unweighted Fibonacci count. Let $\omega$ be a primitive cube root of unity. The exact residue filter and the identity $P_n(1)=F_{n+2}$ give \$a\_n=\textbackslash{}bigl(F\_\{n+2\}+P\_n(\textbackslash{}omega)+P\_n(\textbackslash{}omega\textasciicircum{}2)\textbackslash{}bigr)/3\$, and hence $\frac{a_n}{F_{n+2}}-\frac13=\frac{P_n(\omega)+P_n(\omega^2)}{3F_{n+2}}$. Thus the desired density is isolated from the exact counting problem: it holds precisely when the combined nontrivial contribution is negligible relative to $F_{n+2}$.

The two nontrivial evaluations are complex conjugates, so their sum is $2\operatorname{Re}P_n(\omega)$. In particular, $\left|\frac{a_n}{F_{n+2}}-\frac13\right|\leq\frac{2|P_n(\omega)|}{3F_{n+2}}$. A proof that $|P_n(\omega)|=o(F_{n+2})$ would therefore suffice, although the exact requirement is only $P_n(\omega)+P_n(\omega^2)=o(F_{n+2})$. This distinction separates a convenient bound on an individual complex evaluation from the cancellation actually needed by the counting ratio.

The proposed three-step transfer analysis offers a route to such a bound, but only after the ordered block identities have been justified. For each residue class of the index modulo $3$, the argument would need to identify the growth rates of the nontrivial modes and show that the terms occurring in the corresponding solution grow strictly more slowly than the Fibonacci component. Characteristic-root information must therefore be accompanied by control of the coefficients selected by the initial data; a possible growing mode cannot simply be discarded. The supplied argument establishes neither this comparison nor the required coefficient control, so the proposed block relations cannot yet be used as an asymptotic estimate.

A global recurrence and the limiting density consequently remain distinct proof goals: even a proved annihilating relation would still require an analysis of the solution determined by the counting data. Bounded exact recurrence checks address neither that growth comparison nor the decay of the error displayed above. The supported weighted and gap-removal encodings identify the counting sequence, while the roots-of-unity decomposition identifies the precise additional statement needed for equidistribution: the nontrivial residue contribution must disappear after normalization by the unweighted Fibonacci count.

\section{Verification status and limitations}
\begin{itemize}
\item The general recurrence and limiting density are not proved by the supplied argument summaries. Their unresolved steps belong in the transfer and asymptotic sections rather than being hidden behind finite agreement.
\item The exact value $a_{12}=128$ is supported as a reported finite-count result. This plan neither reruns that computation nor claims independent verification of the saved output.
\item The indexed material gives the initial counts through $a_8$ but does not display a complete initial-value vector through $a_{11}$ for the proposed order-twelve recurrence. The plan must not infer that omitted values are absent from every archived source.
\item The Gaussian-binomial formula has mixed proof-status assessments: its bijection and corrected normalization are supported, while a later review still requests a complete all-index symbolic argument.
\item The evidence index is a bounded selection with omitted notes and fragments. Section assignments identify relevant sources for drafting; they do not certify exhaustive coverage of the archived research.
\end{itemize}
\section{Open obligations}
\begin{itemize}
\item Write and verify the ordered three-step transfer product for every starting phase, without assuming that periodic matrices commute.
\item Establish the claimed block characteristic relations symbolically before using them to derive the candidate scalar recurrence.
\item Check the recurrence's starting index, especially $n=0$, using valid initial data rather than extending the weighted recurrence beyond its stated domain.
\item Supply and justify enough initial values to determine the sequence under the proposed recurrence; for an order-twelve formulation starting at $n=0$, this requires a complete specification through $a_{11}$.
\item Resolve the mixed assessments of the Gaussian-binomial representation by presenting a complete symbolic argument under the corrected convention, including the $k=0$ and endpoint cases.
\item Prove that the combined nontrivial root-of-unity contribution is $o(F_{n+2})$, with a justified growth estimate rather than finite ratio evidence.
\end{itemize}
\section{Conclusion}
The explicit counterexample \$a\_5\textbackslash{}ne a\_4+a\_2\$ rules out the short recurrence $a_{n+3}=a_{n+2}+a_n$. This negative result demonstrates why early pattern agreement cannot substitute for a structural derivation. [C08]

The reported exact finite counts give $a_{12}=128$. This answers the finite evaluation component at the level supported by the supplied computation records, without establishing a general recurrence. [C09]

The order-twelve candidate has bounded computational support, but its transfer-matrix derivation, boundary treatment, and sufficient initial-value specification remain unfinished. It should be presented as a conjectural recurrence, not the requested proved theorem. [C10, C11]

The roots-of-unity decomposition reduces the density question to control of the nontrivial modes. Because that growth estimate is missing, the limit $a_n/F_{n+2}\to1/3$ remains an open obligation within this evidence package. [C12]

\section*{Authorship note}
AI-assisted synthesis from saved research; not independently peer reviewed.
\end{document}
